\scanpage{70}

\textcircled{1}\quad
$k_*\xrightarrow{\varepsilon_*}A_*^*$,
$K_*=(K_n)_{n\geq0}$,
$K_n:A_n^n\text{---}k$, satisf.
$\left\{\begin{array}{l}\text{Stokes}\\ K_0\varepsilon_0=\mathrm{id}\end{array}\right.$

\textcircled{2}\quad Sur un \deleted{\unsure{topos}} espace topologique $X$,
faisceau différentiel \inserted{de $S$-modules}
\deleted{\unsure{$E$}} $\Omega^*$ qui est une \uline{résolution} de $k_X$
(p.ex. complexe de DR sur $X$ variété $C^\infty$, $k=\mathbb R=S$)

\textcircled{3}\quad Un sous-ensemble $1/2$ simplicial $S(X)'$ du complexe singulier de $X$, tel que
$S(X)'\text{---}S(X)$ \deleted{\unsure{soit}} induise un isom sur les \unsure{gr.} d'homologie entière
(p.ex. $X$ est une variété $C^\infty$ et $S(X)'$ est formé des simplexes $C^\infty$)
\inserted{\unsure{et si on remplace} $X$ \unsure{par un ouvert}
(d'une base) et qu'on \unsure{prend dans} $S(X)$ et $S(X)'$
\unsure{les sous-complexes des simplexes} de \uline{support $\subset U$}!
\deleted{\unsure{$J$}}}

\margin{donc, si on regarde le \uline{faisceau} \deleted{\unsure{de}}
de \unsure{coch.} sing. $S^*_{X,k}$ sur $X$, et de même
$S_{X,k}^{*'}$, $S_{X,k}^{*'}$ est comme $S_{X,k}^*$ une rés. de $k$}
\figure{\begin{tikzpicture}\draw[->] (4,0) .. controls (3,.7) and (1,.6) .. (0,.6);\end{tikzpicture}}

\textcircled{4}\quad $\forall\sigma\in S(X)'$, est
\deleted{\unsure{donné un morphisme}}
\inserted{\unsure{un hom. de}} $k$-augm.
\[
\sigma^*:\underset{\substack{U\text{ ouvert de }X\\ U\supset\operatorname{supp}\sigma}}{\varinjlim}
\Omega^*(U)\longrightarrow A_n^*
\qquad(n=\dim\sigma)
\]

On suppose que si $\alpha:\Delta_m\longrightarrow\Delta_n$,
$\sigma\in S'(X)_n$,
\deleted{\unsure{$\sigma:|\Delta_n|\longrightarrow X$}}
on a commutativité dans
\diagram{\begin{tikzcd}
\underset{U\supset\operatorname{supp}\sigma}{\varinjlim}\Omega^*(U)
\arrow[r]\arrow[d,"\text{restriction}"'] & A_n^*\arrow[d,"\unsure{\sigma_A^*|\alpha_A^*}"]\\
\underset{V\supset\operatorname{supp}\alpha^*(\sigma)}{\varinjlim}\Omega^*(V)
\arrow[r] & A_m^*
\end{tikzcd}}
\margin{\uline{NB} $\operatorname{supp}\alpha^*(\sigma)\subset\operatorname{supp}\sigma$}

[\uline{NB} On peut dire que
\[
\sigma\longmapsto\underset{U\supset\operatorname{supp}\sigma}{\varinjlim}\Omega^*(U)
\]
forme un \deleted{\unsure{faisceau}} objet semi-simplicial
$\Omega^*_{S(X)'}$ (en fait un complexe de $S$-modules
$k$-augmenté) sur $S(X)'=\mathcal X$,
et on se donne un hom de celui-ci dans le complexe \unsure{analogue}
$A^*_{\mathcal X}$ sur $\mathcal X=S(X)'$ défini par
$A_*^*\times\mathcal X$.

\margin{\diagram{\begin{tikzcd}
\Delta_n\arrow[dr,"\sigma"'] & \Delta_m\arrow[l,"\alpha"']\arrow[d,"\alpha^*(\sigma)=\sigma\circ\alpha"]\\
& S(X)'=\mathcal X
\end{tikzcd}}}

\uline{NB} finalement \textcircled{1} $(A_*^*,k_*,K_*)$
\inserted{et (4)} ne servent que pour pouvoir définir
\[
\unsure{\varphi:|\leftrightarrow}\quad\Omega^*\longrightarrow S_X^{*'}
\quad\unsure{\text{comp. coch. augm.}}
\]
\margin{69}
